
\chapter{Conclusões e Trabalhos Futuros}
\label{chap:conclusoes-e-trabalhos-futuros}

O presente trabalho teve como objetivo desenvolver a plataforma WebGIS ``Sobral em Mapas'', destinada à organização, ao gerenciamento e à visualização de dados geoespaciais do município de Sobral, no Ceará. Para isso, foram aplicados conceitos de Engenharia de Software, desde o levantamento de requisitos e a modelagem dos dados até a implementação de uma arquitetura desacoplada, utilizando React e Leaflet no \textit{front-end}, Laravel no \textit{back-end} e MySQL na persistência administrativa. Como resultado, foi desenvolvida uma aplicação \textit{web} funcional que reúne ferramentas cartográficas e recursos administrativos para consultar informações territoriais, gerenciar camadas temáticas e controlar o acesso dos usuários.

Entre as principais contribuições, destaca-se a integração de diferentes fontes geoespaciais, inicialmente baseada nos serviços WMS do GeoServer e posteriormente ampliada para arquivos GeoJSON. Essa adaptação reduziu a dependência exclusiva de serviços externos, possibilitando o envio, a pré-visualização e a criação de geometrias diretamente na plataforma. Também foram implementadas ferramentas de pesquisa, medição, exportação e personalização cartográfica, enquanto o uso do Docker no ambiente local e a implantação independente do \textit{front-end} e do \textit{back-end} contribuíram para a organização e a padronização da execução dos componentes. Dessa forma, os resultados demonstram a viabilidade técnica de integrar gerenciamento administrativo e visualização geoespacial em uma plataforma voltada à gestão municipal e à consulta integrada de informações sobre o território.

Além das verificações funcionais, foi realizada uma avaliação exploratória de usabilidade com usuários voluntários, cujos resultados indicaram uma percepção favorável quanto à facilidade de utilização, à satisfação e à utilidade da plataforma. A avaliação também permitiu identificar dificuldades pontuais, especialmente nas operações administrativas de gerenciamento de camadas, contribuindo para orientar ajustes na interface e reconhecer oportunidades de aprimoramento. Embora os resultados tenham sido positivos, o caráter exploratório do estudo e o número limitado de participantes indicam a necessidade de avaliações complementares com diferentes perfis de usuários, incluindo profissionais e gestores públicos.

Apesar dos avanços alcançados, permanecem limitações relacionadas à disponibilidade dos serviços WMS externos, à necessidade de armazenamento persistente dos arquivos GeoJSON no ambiente de hospedagem e ao processamento de conjuntos extensos de dados geográficos. Esses aspectos poderão ser aprimorados por meio de melhorias na infraestrutura, mecanismos de \textit{cache}, simplificação de geometrias e carregamento progressivo de camadas, favorecendo o desempenho e a disponibilidade da plataforma em diferentes condições de utilização.

Como trabalhos futuros, propõe-se ampliar o suporte a formatos geográficos, como KML, e a serviços baseados no padrão WFS, além de incorporar filtros mais detalhados sobre os atributos das camadas e ferramentas adicionais de análise espacial. Também se pretende aprimorar o acompanhamento da origem e da atualização das informações geográficas, favorecendo a organização, a interoperabilidade e a confiabilidade dos dados disponibilizados. Essas melhorias poderão ampliar as possibilidades de utilização da plataforma em atividades de consulta, planejamento e gestão territorial.

Outra possibilidade de evolução consiste no desenvolvimento de uma versão móvel do ``Sobral em Mapas'', com recursos de consulta cartográfica, geolocalização e acesso a informações selecionadas em situações de baixa conectividade. Essa iniciativa poderá começar pela implementação de uma \textit{Progressive Web App} (PWA) ou avançar para uma aplicação desenvolvida com React Native, aproveitando os conhecimentos e a lógica de integração já utilizados no projeto. Também se considera a integração de um assistente virtual (\textit{chatbot}), utilizando tecnologias como o \textit{framework} Rasa para interpretar perguntas em linguagem natural, orientar os usuários e facilitar o acesso às ferramentas e camadas do mapa. Ademais, pretende-se dar continuidade às melhorias de acessibilidade, responsividade e personalização da interface, considerando as necessidades identificadas em futuras avaliações.

Conclui-se, portanto, que o ``Sobral em Mapas'' alcançou seu objetivo principal ao reunir visualização cartográfica e gerenciamento administrativo em uma plataforma WebGIS funcional. A arquitetura desacoplada e o suporte a diferentes fontes geoespaciais estabeleceram uma base tecnológica que favorece a manutenção e a evolução do sistema, enquanto a avaliação exploratória apresentou indícios favoráveis de sua aceitação pelos participantes. Assim, a plataforma demonstra potencial para contribuir com a organização e a disponibilização de informações territoriais, oferecendo recursos que podem apoiar consultas públicas e atividades de planejamento e gestão do município de Sobral.
